% Class-compatibility preamble retained for a reader excerpt of
%   Dave Witte Morris, "A Cohomological Proof that Real Representations of
%   Semisimple Lie Algebras Have Q-Forms", SIGMA 11 (2015) 034,
%   DOI 10.3842/SIGMA.2015.034, CC BY 4.0.
%
% The source pins the publisher class sigma, which is not part of the delivery
% baseline, so this file supplies the amsart baseline and then copies the author's
% own declaration block mechanically and unchanged:
%   - pinned-original-source.tex lines 11-31: the author's DeclareMathOperator
%     list, the equation numbering and the author's newtheorem declarations;
%   - pinned-original-source.tex lines 35-43, which the author writes immediately
%     after \begin{document} in the publisher layout: the author's own \Res,
%     \ResOne and \Lie macros and \allowdisplaybreaks.
% The publisher header commands of the sigma class (\FirstPageHeading,
% \ArticleName and the like) and the macro \PaperNumber they require are not
% reproduced, because this excerpt carries neither the journal front matter nor
% anything else that uses them.
% No mathematical macro, statement, symbol or assumption is introduced, changed
% or removed by this file.  The author's full DeclareMathOperator list of pinned
% source lines 11-19 is carried, including GL, SL, SO and SU: these four are not
% supplied by the delivery baseline, so leaving them out makes the author's own
% notation (for instance GL_n(Q) in the corollary on lattices) undefined.
\documentclass[10pt,a4paper]{amsart}
\usepackage{amssymb,amsmath,amsthm}
\usepackage{enumerate}
\usepackage{hyperref}

\DeclareMathOperator{\Cliff}{Clif\/f}
\DeclareMathOperator{\Gal}{Gal}
\DeclareMathOperator{\GL}{GL}
\DeclareMathOperator{\Inn}{Inn}
\DeclareMathOperator{\Mat}{Mat}
\DeclareMathOperator{\Out}{Out}
\DeclareMathOperator{\SL}{SL}
\DeclareMathOperator{\SO}{SO}
\DeclareMathOperator{\SU}{\mathbf{SU}}
\numberwithin{equation}{section}

\newtheorem{prop}{Proposition}[section]
\newtheorem{cor}[prop]{Corollary}
\newtheorem{lem}[prop]{Lemma}
{\theoremstyle{definition}
\newtheorem{defn}[prop]{Definition}
\newtheorem{notation}[prop]{Notation}
\newtheorem{rem}[prop]{Remark}
\newtheorem{ques}[prop]{Question}
\newtheorem{eg}[prop]{Example}
\newtheorem{erratum}[prop]{Erratum}}

\newcommand{\Res}[1]{\mathop{\mathbf{R}_{#1}}}
\newcommand{\ResOne}[1]{\mathop{\mathbf{R}^{(1)}_{#1}}}
\newcommand{\Lie}[1]{\mathfrak{\lowercase{#1}}}
\allowdisplaybreaks
\usepackage{cleveref}
\crefname{originallemma}{original lemma}{original lemmas}
\Crefname{originallemma}{Original lemma}{Original lemmas}
\makeatletter
\newlabel{SameFixed}{{2.3}{source}{}{sourcepointer-1}{}}
\newlabel{SameFixed@cref}{{[originallemma][0][]2.3}{[1][0][]source}{}{}{}}
\makeatother
\crefname{originallemma}{original lemma}{original lemmas}
\Crefname{originallemma}{Original lemma}{Original lemmas}
\makeatletter
\newlabel{RunivIff}{{2.1}{source}{}{sourcepointer-2}{}}
\newlabel{RunivIff@cref}{{[originallemma][0][]2.1}{[1][0][]source}{}{}{}}
\makeatother
\crefname{originaldefinition}{original definition}{original definitions}
\Crefname{originaldefinition}{Original definition}{Original definitions}
\makeatletter
\newlabel{RunivIff-D}{{4}{source}{}{sourcepointer-3}{}}
\newlabel{RunivIff-D@cref}{{[originaldefinition][0][]4}{[1][0][]source}{}{}{}}
\makeatother
\title{R-universality of the special unitary Lie algebra of a Hermitian form over a quaternion division algebra}
\author{Dave Witte Morris}
\date{doi:10.3842/SIGMA.2015.034; CC BY 4.0}
\begin{document}
\maketitle
\noindent\textit{Editorial note.} All delivered mathematical text is the author's own, copied from the pinned publisher source sources/algebra-fl02-a-cohomological-proof-that-real-re/pinned-original-source.tex; no statement or proof step was rewritten, reordered or removed.\par\medskip
\noindent\textit{Editorial note.} Preamble provenance: the source pins the publisher class sigma, which is not part of the delivery baseline, so this card uses sources/algebra-fl02-a-cohomological-proof-that-real-re/retained-preamble.tex, which supplies the amsart baseline, copies the author's DeclareMathOperator list, equation numbering and newtheorem declarations from pinned source lines 11-31, and copies the author's own \textbackslash{}Res, \textbackslash{}ResOne and \textbackslash{}Lie macros from pinned source lines 35-43. No macro was invented and no mathematical symbol was changed.\par\medskip
\noindent\textit{Editorial note.} Bibliography provenance: the citations of the counted proof are resolved with sources/algebra-fl02-a-cohomological-proof-that-real-re/bibliography.tex, an excerpt of the author's own thebibliography environment (pinned source lines 954-1016) copied character for character, all entries kept so that every citation of a delivered excerpt resolves; no entry was rewritten, reordered, merged, split or invented and no citation key was changed.\par\medskip
\noindent\textit{Editorial note.} External references are navigation pointers only; every deliberately kept step of the author's argument stays inside the counted spans. Their printed numbers are the author's own numbering under the author's newtheorem declarations, computed mechanically by the lane's scanner.\par\medskip
\noindent\textit{Allowed standard prerequisites (editorial).}\begin{itemize}
\item The Tits algebra of the natural representation of the unitary algebra is the division algebra D. Used inside the counted proof at source lines 585-586 and again at source lines 591-592, attributed by the author to page 378 of the cited book KnusEtAl-BookOfInvols. Identifies the Tits algebra whose base change decides universality.
\item The minimal-extension lemma for the star-action. The author's lemma labelled SameFixed, used inside the counted proof at source line 579. Makes the quadratic extension L imaginary because the algebra is outer.
\end{itemize}\medskip
\ifdefined\hypertarget\hypertarget{sourcepointer-1}{}\fi
\noindent\textit{Source reference.} Original lemma 2.3: sources/algebra-fl02-a-cohomological-proof-that-real-re/pinned-original-source.tex, source lines 471--471. Navigation pointer only: the counted proof uses this lemma once to conclude that the extension is imaginary; its content is declared as a named standard prerequisite and it has its own card in this lane.\par\medskip
\ifdefined\hypertarget\hypertarget{sourcepointer-2}{}\fi
\noindent\textit{Source reference.} Original lemma 2.1: sources/algebra-fl02-a-cohomological-proof-that-real-re/pinned-original-source.tex, source lines 416--416. Navigation pointer only: the counted proof invokes the author's equivalence for its closing steps.\par\medskip
\ifdefined\hypertarget\hypertarget{sourcepointer-3}{}\fi
\noindent\textit{Source reference.} Original definition 4: sources/algebra-fl02-a-cohomological-proof-that-real-re/pinned-original-source.tex, source lines 432--432. Navigation pointer only: the labelled item of that equivalence used at source line 588.\par\medskip
\begin{prop}
Let $\Lie G = \Lie{SU}_n(B; D, \tau)$, where
\begin{itemize}\itemsep=0pt
\item $D$ is a~division algebra that is central over a~quadratic extension~$L$ of~$\mathbb{Q}$,
\item $\tau$ is an anti-involution of~$D$ that is nontrivial on~$L$, and
\item $B$ is an invertible~$\tau$-Hermitian matrix in $\Mat_{n\times n}(D)$.
\end{itemize}
Then $\Lie G$ is $\mathbb{R}$-universal if and only if $D = L$ is an imaginary quadratic extension
of~$\mathbb{Q}$, and either~$n$~is odd or $(-1)^{n/2} \det B$ is either negative or the norm of some element of~$L$.
\end{prop}

\begin{proof}
We prove only ($\Rightarrow$), but the argument is reversible.
Since $\Lie G$ is outer, Lemma~\ref{SameFixed} implies that $\Lie G \otimes_\mathbb{Q} \mathbb{R}$ is also outer, and the
quadratic extension~$L$ is imaginary.

Let~$d$ be the degree of~$D$ over~$L$, and let~$\lambda$ be the highest weight of the standard representation of
$\Lie{SL}_{dn}(\mathbb{C})$, so $\{\lambda^i\}_{i=1}^{dn}$ is a~set of representatives of~$Z^*$.

The Tits algebra of the natural representation $\rho \colon \Lie G \hookrightarrow \Mat_{n \times n}(D)$
is~$D$~\cite[p.~378]{KnusEtAl-BookOfInvols}.
Since~$D$ splits over~$\mathbb{R}$ (recall that the center~$L$ is an imaginary extension),
we have $D_{\Lie G, \mathbb{R}}(\lambda) = \mathbb{C}$, so we conclude from Lemma~\ref{RunivIff}(\ref{RunivIff-D})
that $D = L$.

If~$\lambda$ is any weight that is not f\/ixed by the $*$-action, then $D_{\Lie G, \mathbb{Q}}(\lambda) = L$ and $D_{\Lie
G, \mathbb{R}}(\lambda) = \mathbb{C}$~\cite[p.~378]{KnusEtAl-BookOfInvols}, so $D_{\Lie G, \mathbb{Q}}(\lambda)
\otimes_\mathbb{Q} \mathbb{R} = D_{\Lie G, \mathbb{R}}(\lambda)$, as specif\/ied in Lemma~\ref{RunivIff}(\ref{RunivIff-D}).
Hence, such weights (which are all of the weights when~$n$ is odd) do not place any further restriction on~$\Lie G$.

Suppose~$n$ is even.
Any weight that is f\/ixed by the $*$-action (and is not in the root lattice) is congruent to $\lambda^{n/2}$, modulo the
root lattice.
Write $L = \mathbb{Q} \bigl[\sqrt{a} \bigr]$ and let $b = (-1)^{n/2} \det B$.
Then, for~$F$ either~$\mathbb{Q}$ or~$\mathbb{R}$, $D_{\Lie G, F}(\lambda^{n/2})$ is Brauer equivalent to the quaternion
algebra $(a,b)_F$~\cite[p.~378 and Corollary~10.35 on p.~131)]{KnusEtAl-BookOfInvols}.
This is trivial in the Brauer group if and only if it is split, which means that~$b$ is the norm of some element of $F
\bigl[\sqrt{a} \bigr]$.

Since $\Lie G$ is $\mathbb{R}$-universal, we see from Lemma~\ref{RunivIff}(\ref{RunivIff-D}) that either $D_{\Lie G,
\mathbb{R}}(\lambda^{n/2}) = \mathbb{H}$ (which, by the preceding paragraph, means that~$b$ is not a~norm in $\mathbb{R}
\bigl[\sqrt{a} \bigr] = \mathbb{C}$, so $b < 0$) or $D_{\Lie G, \mathbb{Q}}(\lambda^{n/2}) = \mathbb{Q}$ (which means
that~$b$~is a~norm in $\mathbb{Q} \bigl[\sqrt{a} \bigr] = L$).
\end{proof}

%% Bibliography excerpt retained for reader excerpts of
%   Dave Witte Morris, "A Cohomological Proof that Real Representations of Semisimple Lie Algebras Have Q-Forms",
%   SIGMA 11 (2015) 034, DOI 10.3842/SIGMA.2015.034.
%%
%% Every entry below is copied character for character from the author's own
% thebibliography environment of the pinned publisher source (lines 954-1016).
%% No entry was rewritten, reordered, merged, split or invented, and no citation
% key was changed.  All 12 entries are kept so that any citation of a delivered
%% excerpt resolves without guessing an entry.
\begin{thebibliography}{99}

\footnotesize\itemsep=0pt

\bibitem{DemazureGrothendieck-SGA3}
Demazure M., Grothendieck A., Sch\'emas en groups.~III. Structure des sch\'emas
  en groupes reductifs, S\'eminaire de g\'eom\'etrie alg\'ebrique du {B}ois
  {M}arie, 1962/1964, available at
  \url{http://library.msri.org/books/sga/sga/pdf/sga3-3.pdf}.

\bibitem{Harder-BerichtNeuere}
Harder G., Bericht \"uber neuere {R}esultate der {G}aloiskohomologie
  halbeinfacher {G}ruppen, \textit{Jber. Deutsch. Math.-Verein.} \textbf{70}
  (1968), 182--216, available at
  \url{http://www.digizeitschriften.de/en/dms/img/?PPN=PPN37721857X_0070&DMDID=dmdlog17}.

\bibitem{Kneser-LectGalCohoClassical}
Kneser M., Lectures on {G}alois cohomology of classical groups, \textit{Tata
  Institute of Fundamental Research Lectures on Mathematics}, Vol.~47, Tata
  Institute of Fundamental Research, Bombay, 1969, available at
  \url{http://www.math.tifr.res.in/~publ/ln/tifr47.pdf}.

\bibitem{KnusEtAl-BookOfInvols}
Knus M.A., Merkurjev A., Rost M., Tignol J.P., The book of involutions,
  \textit{American Mathematical Society Colloquium Publications}, Vol.~44,
  Amer. Math. Soc., Providence, RI, 1998.

\bibitem{Lam-IntroQuadForms}
Lam T.Y., Introduction to quadratic forms over f\/ields, \textit{Graduate Studies
  in Mathematics}, Vol.~67, Amer. Math. Soc., Providence, RI, 2005.

\bibitem{Morris-RRepsQForms}
Morris D., Real representations of semisimple {L}ie algebras have {${\mathbb
  Q}$}-forms, in Algebraic Groups and Arithmetic, Tata Institute of Fundamental
  Research, Mumbai, 2004, 469--490.

\bibitem{PlatonovRapinchukBook}
Platonov V., Rapinchuk A., Algebraic groups and number theory, \textit{Pure and
  Applied Mathematics}, Vol.~139, Academic Press, Inc., Boston, MA, 1994.

\bibitem{PrasadRapinchuk-Prescribed}
Prasad G., Rapinchuk A., On the existence of isotropic forms of semi-simple
  algebraic groups over number f\/ields with prescribed local behavior,
  \href{http://dx.doi.org/10.1016/j.aim.2006.01.001}{\textit{Adv. Math.}} \textbf{207} (2006), 646--660.

\bibitem{Raghunathan-ArithLattsSSGrps}
Raghunathan M.S., Arithmetic lattices in semisimple groups, \href{http://dx.doi.org/10.1007/BF02967980}{\textit{Proc.
  Indian Acad. Sci. Math. Sci.}} \textbf{91} (1982), 133--138.

\bibitem{Tits-Classification}
Tits J., Classif\/ication of algebraic semisimple groups, in Algebraic {G}roups
  and {D}iscontinuous {S}ubgroups ({P}roc. {S}ympos. {P}ure {M}ath., {B}oulder,
  {C}olo., 1965), Amer. Math. Soc., Providence, R.I., 1966, 33--62.

\bibitem{Tits-RepLinIrred}
Tits J., Repr\'esentations lin\'eaires irr\'eductibles d'un groupe r\'eductif
  sur un corps quelconque, \href{http://dx.doi.org/10.1515/crll.1971.247.196}{\textit{J.~Reine Angew. Math.}} \textbf{247} (1971),
  196--220.

\bibitem{Tits-StronglyInner}
Tits J., Strongly inner anisotropic forms of simple algebraic groups,
  \href{http://dx.doi.org/10.1016/0021-8693(90)90201-X}{\textit{J.~Algebra}} \textbf{131} (1990), 648--677.


\end{thebibliography}
\end{document}
